% Part: incompleteness
% Chapter: arithmetization-syntax
% Section: proofs-in-nd

\documentclass[../../../include/open-logic-section]{subfiles}

\begin{document}

\olfileid{inc}{art}{pnd}
\olsection{\usetoken{P}{derivation} ing Deduksi Natural}

\begin{explain}
Supaya bisa ngaritmetisasi !!{derivation}s, kita kudu makili
!!{derivation}s minangka wilangan. Amarga !!{derivation}s iku wit
saka !!{formula}s kang saben inferensine nggawa siji utawa loro
label, cara kang paling cetha yaiku representasi rekursif: kita
makili !!{derivation} minangka tupla kang komponene yaiku cacah
sub-!!{derivation}s langsung kang tumuju marang premis inferensi
pungkasan, representasi sub-!!{derivation}s mau, !!{formula}
pungkasan, label pambebasan saka inferensi pungkasan, lan wilangan
kang nuduhake jinis inferensi pungkasan.
\end{explain}

\begin{defn}
  Yen $\delta$ iku !!a{derivation} ing deduksi natural, $\Gn{\delta}$
  ditemtokake kanthi induktif kaya mangkene:
  \begin{enumerate}
  \item 
    Yen $\delta$ mung dumadi saka asumsi~$!A$, $\Gn{\delta}$
    yaiku $\tuple{0, \Gn{!A}, n}$. Wilangan~$n$ yaiku~$0$
    yen asumsi iku !!{undischarged}, lan dadi label wilangane
    yen ora mangkono.
  \item 
    Yen $\delta$ dipungkasi inferensi kanthi nol, siji, loro,
    utawa telu premis, $\Gn{\delta}$ yaiku
    \begin{align*}
      & \tuple{0, \Gn{!A}, n, k}, \\
      & \tuple{1, \Gn{\delta_1}, \Gn{!A}, n, k}, \\
      & \tuple{2, \Gn{\delta_1}, \Gn{\delta_2}, \Gn{!A}, n, k}, \text{ utawa}\\
      & \tuple{3, \Gn{\delta_1}, \Gn{\delta_2}, \Gn{\delta_3}, \Gn{!A}, n,
        k},
    \end{align*}
    miturut urutane. Ing kene $\delta_1$, $\delta_2$, $\delta_3$
    iku sub-!!{derivation}s kang pungkasane premis inferensi
    pungkasan ing $\delta$; $!A$ iku konklusi inferensi pungkasan
    ing~$\delta$; $n$ iku label pambebasan inferensi pungkasan
    ($0$ yen inferensi iku ora mbebasake asumsi); lan $k$~diwenehi
    dening tabel iki miturut aturan kang digunakake ing inferensi
    pungkasan.
  
    \begin{tabular}{lccccccc}
      \text{Aturan:} & \Intro{\land} & \Elim{\land} & \Intro{\lor} & \Elim{\lor} \\
      $k$: & 1 & 2 & 3 & 4  \\[1.5ex]
      \text{Aturan:} &  \Intro{\lif} & \Elim{\lif} & \Intro{\lnot} & \Elim{\lnot} \\
      $k$: & 5 & 6 & 7 & 8 \\[1.5ex]
      \text{Aturan:}  &  \FalseInt & \FalseCl & \Intro{\lforall} & \Elim{\lforall} \\
      $k$: & 9 & 10 & 11 & 12 \\[1.5ex]
      \text{Aturan:} & \Intro{\lexists} & \Elim{\lexists} & \Intro{\eq} & \Elim{\eq} \\
      $k$: & 13 & 14 & 15 & 16 
    \end{tabular}
  \end{enumerate}
\end{defn}

\begin{ex}
  Gatekna !!{derivation} kang prasaja iki
  \begin{prooftree}
    \AxiomC{$\Discharge{!A \land !B}{1}$}
    \RightLabel{\Elim{\land}}
    \UnaryInfC{$!A$}
    \DischargeRule{\Intro{\lif}}{1}
    \UnaryInfC{$(!A \land !B) \lif !A$}
  \end{prooftree}
  Wilangan G\"odel saka asumsi iku $d_0 = \tuple{0,
    \Gn{!A \land !B}, 1}$. Wilangan G\"odel saka !!{derivation}
  kang pungkasane konklusi~$\Elim{\land}$ iku $d_1 = \tuple{1,
     d_0, \Gn{!A}, 0, 2}$ ($1$ amarga $\Elim{\land}$ nduweni siji
  premis, banjur wilangan G\"odel saka konklusi~$!A$, $0$ amarga
  ora ana asumsi kang dibebasake, lan $2$ iku wilangan kang ngodeake
  $\Elim{\land}$). Mula wilangan G\"odel saka sakabehe
  !!{derivation} yaiku
  $\tuple{1, d_1, \Gn{((!A \land !B) \lif
      !A)}, 1, 5}$, yaiku
  \[
  \tuple{1, \tuple{1, \tuple{0, \Gn{(!A \land !B)}, 1}, \Gn{!A}, 0, 2},
    \Gn{((!A \land !B) \lif !A)}, 1, 5}.
  \]
\end{ex}

\begin{explain}
Sawise nemtokake representasi !!{derivation}s, kita uga kudu
nuduhake yen wilangan G\"odel saka !!{derivation}s iki bisa
diolah kanthi rekursif primitif, lan sifat sarta relasi pokoke
bisa dinyatakake. Sawetara operasi prasaja: upamane, yen diwenehi
wilangan G\"odel~$d$ saka !!a{derivation},
$\fn{EndFmla}(d) = (d)_{(d)_0+1}$ menehi wilangan G\"odel saka
!!{formula} pungkasan, $\fn{DischargeLabel}(d) = (d)_{(d)_0 + 2}$
menehi label pambebasan, lan $\fn{LastRule}(d) = (d)_{(d)_0 + 3}$
menehi wilangan kang nuduhake jinis inferensi pungkasan.
Sawetara sifat liyane luwih angel. Kita bakal menehi sketsa carane.
Tujuane yaiku nuduhake yen relasi ``$\delta$ iku !!a{derivation}
saka~$!A$ adhedhasar~$\Gamma$'' iku relasi rekursif primitif
saka wilangan G\"odel $\delta$ lan~$!A$.
\end{explain}

\begin{prop}
  Relasi-relasi iki rekursif primitif:
  \begin{enumerate}
  \item $!A$ ana minangka asumsi ing~$\delta$ kanthi label~$n$.
  \item Kabeh asumsi ing $\delta$ kanthi label~$n$ awujud~$!A$
    (yaiku, kita bisa !!{discharge} asumsi~$!A$ nganggo label~$n$
    ing~$\delta$).
  \end{enumerate}
\end{prop}

\begin{proof}
  Kita kudu nuduhake yen relasi kang cocog ing antarane wilangan
  G\"odel saka !!{formula}s lan wilangan G\"odel saka
  !!{derivation}s iku rekursif primitif.
  \begin{enumerate}
  \item Kita arep nuduhake yen $\fn{Assum}(x, d, n)$ iku rekursif
    primitif. Relasi iki bener yen $x$ iku wilangan G\"odel saka
    asumsi ing !!{derivation} kang wilangan G\"odele~$d$
    kanthi label~$n$. Iki lumaku yen !!{derivation} kanthi wilangan
    G\"odel~$\tuple{0, x, n}$ iku sub-!!{derivation} saka~$d$.
    Gatekna yen cara kita ngodeake !!{derivation}s iku kasus khusus
    saka pangodean wit kang dikenalake ing
    \olref[cmp][rec][tre]{sec}. Mula fungsi rekursif primitif
    $\fn{SubtreeSeq}(d)$ menehi urutan wilangan G\"odel saka kabeh
    sub-!!{derivation}s saka~$d$ (kanthi dawa kang uga rekursif primitif).
    Dadi kita bisa netepake
    \[
    \fn{Assum}(x, d, n) \defiff \bexists{i<\len{\fn{SubtreeSeq}(d)}}{(\fn{SubtreeSeq}(d))_i =
      \tuple{0, x, n}}.
    \]
Cathetan koreksi sumber SOL6-F083: SubtreeSeq ing sumber bisa ngemot pangulangan. Kanggo wit siji anak kang dadi godhong, dawa dhaptar sawise d langkah yaiku d tambah siji, mula klaim kabeh dawane paling akeh d ora bener. Assum saiki nggoleki nganti dawa dhaptar kang nyata, yaiku wates rekursif primitif.
  \item Kita arep nuduhake yen $\fn{Discharge}(x, d, n)$, kang
    bener yen kabeh asumsi kanthi label~$n$ ing !!{derivation}
    kanthi wilangan G\"odel~$d$ iku !!{formula} kanthi wilangan
    G\"odel~$x$. Relasi iki bener yen lan mung yen
    $\bforall{y<d}{(\fn{Assum}(y, d, n) \lif y
      = x)}$.
  \end{enumerate}
\end{proof}

\begin{prop}
  \ollabel{prop:followsby}
  Sifat $\fn{Correct}(d)$, kang bener yen lan mung yen inferensi
  pungkasan ing !!{derivation}~$\delta$ kanthi wilangan G\"odel~$d$
  iku bener, yaiku rekursif primitif.
\end{prop}

\begin{proof}
  Ing kene kita kudu nuduhake yen kanggo saben aturan inferensi~$R$,
  relasi $\fn{FollowsBy}_R(d)$ iku rekursif primitif.
  $\fn{FollowsBy}_R(d)$ bener yen lan mung yen $d$ iku wilangan
  G\"odel saka !!{derivation}~$\delta$, lan !!{formula} pungkasane
  $\delta$ diduweni kanthi ngetrapake~$R$ kanthi bener marang
  sub-!!{derivation}s langsung saka~$\delta$.

  Kasus kang prasaja yaiku aturan \Intro{\land}. Yen $\delta$
  dipungkasi inferensi $\Intro{\land}$ kang bener, wujude kaya iki:
  \begin{prooftree}
    \AxiomC{}
    \RightLabel{$\delta_1$}
    \DeduceC{$!A$}
    
    \AxiomC{}
    \RightLabel{$\delta_2$}
    \DeduceC{$!B$}

    \RightLabel{\Intro\land}
    \BinaryInfC{$!A \land !B$}
  \end{prooftree}
  Mula wilangan G\"odel~$d$ saka~$\delta$ yaiku
  $\tuple{2, d_1, d_2, \Gn{(!A \land !B)}, 0, k}$,
  ing ngendi $\fn{EndFmla}(d_1) = \Gn{!A}$,
  $\fn{EndFmla}(d_2) = \Gn{!B}$, $n=0$, lan $k=1$.
  Dadi $\fn{FollowsBy}_{\Intro\land}(d)$ bisa ditetepake minangka
  \begin{multline*}
    (d)_0 = 2 \land \fn{DischargeLabel}(d) = 0 \land \fn{LastRule}(d) = 1 \land {}\\
    \fn{EndFmla}(d) = {}\\
    \Gn{(} \concat \fn{EndFmla}((d)_1) \concat \Gn{\land}
    \concat \fn{EndFmla}((d)_2) \concat \Gn{)}.
  \end{multline*}

  Tuladha prasaja liyane yaiku aturan $\Intro\eq$. Aturan iki
  ora nduweni premis, mula $(d)_0 = 0$, kaya asumsi. Aturan iki
  uga ora nduweni label pambebasan, yaiku $n=0$. Nanging $!A$
  kudu awujud $\eq[t][t]$ kanggo suku katutup~$t$.
  Ing kene definisi rekursif primitife yaiku
  \begin{multline*}
    (d)_0 = 0 \land \fn{DischargeLabel}(d) = 0 \land {}\\
    \bexists{t<d}{(\fn{ClTerm}(t) \land 
      \fn{EndFmla}(d) = {}\\
      \Gn{{\eq}(} \concat t \concat \Gn{,} \concat t \concat \Gn{)})}.
  \end{multline*}

  Kanggo tuladha kang luwih rumit,
  $\fn{FollowsBy}_{\Intro{\lif}}(d)$ bener yen lan mung yen
  !!{formula} pungkasane~$\delta$ awujud $(!A \lif !B)$,
  !!{formula} pungkasane $\delta_1$ yaiku~$!B$, lan nalika ana
  pambebasan, saben asumsi ing~$\delta$ kanthi label~$n$
  awujud~$!A$. Yen label pambebasan nol, ora ana asumsi kang
  dibebasake. Iki bisa
  dinyatakake kanthi rekursif primitif dening
  \begin{multline*}
    (d)_0 = 1 \land {}\\
    \bexists{a<d}{
      \fn{Frm}(a) \land (\fn{DischargeLabel}(d) = 0 \lor
       \fn{Discharge}(a, (d)_1, \fn{DischargeLabel}(d))) \land {}\\
      \fn{EndFmla}(d) = (\Gn{(} \concat a \concat \Gn{\lif}
      \concat \fn{EndFmla}((d)_1) \concat \Gn{)})}
  \end{multline*}
Cathetan koreksi sumber OLPL-179: label pambebasan nol kudu ngidinake introduksi implikasi tanpa asumsi kang dibebasake; rumus Jawa wis nglebokake kasus iki. Klaim lawas ngenani ikatan a ditarik: nutup argumen TeX ora nutup cakupan kuantor kang kacithak.
  (Anggepen $a$ minangka wilangan G\"odel saka~$!A$).

  Kanggo tuladha liyane, gatekna \Intro{\lexists}. Ing kene
  inferensi pungkasan ing~$\delta$ bener yen lan mung yen ana
  !!a{formula}~$!A$, suku katutup~$t$, lan !!a{variable}~$x$
  saengga $\Subst{!A}{t}{x}$ iku !!{formula} pungkasane
  !!{derivation}~$\delta_1$ lan $\lexists[x][!A]$ iku konklusi
  inferensi pungkasan. Mula
  $\fn{FollowsBy}_{\Intro{\lexists}}(d)$ bener yen lan mung yen
  \begin{multline*}
    (d)_0 = 1 \land \fn{DischargeLabel}(d) = 0 \land {} \\
    \bexists{a < d}{\bexists{x<d}{\bexists{t<d}{
          (\fn{Frm}(a) \land \fn{ClTerm}(t) \land \fn{Var}(x)  \land {}}}}\\
    \fn{Subst}(a,t,x) = \fn{EndFmla}((d)_1) \land
    \fn{EndFmla}(d) = (\Gn{\lexists} \concat x \concat a)).
  \end{multline*}

  Sabanjure kita netepake $\fn{Correct}(d)$ minangka
  \begin{multline*}
    \fn{Seq}(d) \land {}\\
    [((d)_0=0 \land \len{d}=3) \lor
    ((d)_0\le3 \land \len{d}=(d)_0+4)] \land {}\\
    \fn{Sent}(\fn{EndFmla}(d)) \land [\\
    (\fn{LastRule}(d) = 1 \land
    \fn{FollowsBy}_{\Intro\land}(d)) \lor \dots \lor {}\\
    (\fn{LastRule}(d) = 16 \land \fn{FollowsBy}_{\Elim\eq}(d)) \lor {}\\
    \bexists{n<d}{\bexists{x<d}{(d = \tuple{0, x, n})}}].
  \end{multline*}
Cathetan koreksi sumber OLPL-176: syarat Sent kudu nyakup kabeh alternatif ing Correct, kalebu kasus asumsi; kurung siku ing rumus Jawa wis njlentrehake cakupan iku.
  Tes Sent mesthekake yen !!{formula} pungkasane~$d$ iku
  ukara. Baris pungkasan nyakup kasus nalika $d$ mung asumsi.
Cathetan koreksi sumber SOL6-F079: Seq nguji kode urutan sah: kode kosong nol utawa siji, utawa produk pangkat positif prima wiwit kapisan tanpa bolong. Tes iki rekursif primitif. Dawa tupla pas mbedakake asumsi lan aturan kanthi cacah premis nol nganti telu.
\end{proof}

\begin{prob}
  Tetepna sifat-sifat iki kaya ing
  \olref[inc][art][pnd]{prop:followsby}:
  \begin{enumerate}
  \item $\fn{FollowsBy}_{\Elim{\lif}}(d)$,
  \item $\fn{FollowsBy}_{\Elim{\eq}}(d)$,
  \item $\fn{FollowsBy}_{\Elim{\lor}}(d)$,
  \item $\fn{FollowsBy}_{\Intro{\lforall}}(d)$.
  \end{enumerate}
  Kanggo sifat pungkasan, kowé uga kudu nuduhake yen kanthi
  rekursif primitif bisa dipriksa apa inferensi pungkasan saka
  !!{derivation} kanthi wilangan G\"odel~$d$ nyukupi syarat
  eigenvariabel, yaiku eigenvariabel~$a$ saka inferensi
  $\Intro{\lforall}$ ora ana ing !!{formula} pungkasane~$d$
  utawa ing asumsi kabuka saka~$d$. Kanggo iki, kowé bisa
  nggunakake predikat rekursif primitif $\fn{OpenAssum}$ saka
  \olref[inc][art][pnd]{prop:openassum}.
\end{prob}

\begin{prop}
  \ollabel{prop:deriv}
  Relasi $\fn{Deriv}(d)$, kang bener yen $d$ iku wilangan G\"odel
  saka !!{derivation}~$\delta$ kang bener, iku rekursif primitif.
\end{prop}

\begin{proof}
  !!^a{derivation}~$\delta$ iku bener yen saben inferensine
  ngetrapake aturan kanthi bener, yaiku yen saben
  sub-!!{derivation}s-e dipungkasi inferensi kang bener.
  Mula $\fn{Deriv}(d)$ bener yen lan mung yen
  \[
  \bforall{i<\len{\fn{SubtreeSeq}(d)}}{\fn{Correct}((\fn{SubtreeSeq}(d))_i)}
  \]
\end{proof}

\begin{prop}
  \ollabel{prop:openassum} Relasi $\fn{OpenAssum}(z, d)$,
  kang bener yen $z$ iku wilangan G\"odel saka asumsi~$!A$
  kang !!a{undischarged} ing !!{derivation}~$\delta$ kanthi
  wilangan G\"odel~$d$, iku rekursif primitif.
\end{prop}

\begin{proof}
  Sawijining panggonan asumsi iku !!{discharged} yen ana kanthi
  label~$n$ ing sub-!!{derivation} saka~$\delta$ kang dipungkasi
  aturan kanthi label pambebasan~$n$. Mula $!A$ iku asumsi kang
  !!a{undischarged} ing~$\delta$ yen paling ora ana siji
  panggonane kang ora !!{discharged} ing~$\delta$. Kita kudu
  ngati-ati: $\delta$ bisa ngemot panggonan~$!A$ kang
  !!{discharged} lan kang !!{undischarged} bebarengan.

  Gatekna urutan $\delta_0$, \dots, $\delta_k$ kanthi $\delta_0 =
  \delta$, $\delta_k$ iku asumsi $\Discharge{!A}{n}$
  (kanggo sawijining~$n$), lan $\delta_{i+1}$ iku sub-!!{derivation}
  langsung saka~$\delta_i$. Yen ana urutan kaya mangkono kang
  ora ana $\delta_i$ dipungkasi inferensi kanthi label
  pambebasan~$n$ kanggo label positif, utawa yen label mau nol, mula $!A$ iku asumsi kang !!a{undischarged}
  ing~$\delta$.

  Fungsi rekursif primitif $\fn{SubtreeSeq}(d)$ menehi urutan
  wilangan G\"odel saka kabeh sub-!!{derivation}s saka~$\delta$.
  Saben urutan wilangan G\"odel saka sub-!!{derivation}s
  saka~$\delta$ nduweni unsur kang uga ana ing urutan mau.
  Syarat saben unsur iki rekursif primitif: relasi kang ing kene
  dijenengi $\fn{Subseq}(s, s')$ bener yen lan mung yen
  $\bforall{i<\len{s}}{\lexists[j<\len{s'}][(s)_i = (s')_j]}$.
Cathetan koreksi sumber OLPL-180: predikat kang dijenengi Subseq ing rumus iki mung mriksa yen saben unsur urutan kapisan ana ing urutan kapindho; rumus iki ora mriksa urutan posisi kang dibutuhake subsekuen sejati. Bukti jalur sabanjure nganggo syarat anak langsung kanthi kapisah.
  Dadi sub-!!{derivation} langsung uga relasi kaya mangkono:
  $\fn{Subderiv}(d, d')$ bener yen lan mung yen
  $\bexists{j<(d')_0}{d = (d')_{j+1}}$. Mula
  $\fn{OpenAssum}(z, d)$ bisa ditetepake dening
  \begin{multline*}
    \bexists{s\le\fn{SubtreeSeq}(d)}{
      \fn{Seq}(s) \land \len{s}>0 \land {}\\
      \fn{Subseq}(s, \fn{SubtreeSeq}(d))
      \land (s)_0 = d \land {}\\
      \bexists{n<d}{(s)_{\len{s} \tsub 1} = \tuple{0, z, n} \land {}\\
        \bforall{i<(\len{s} \tsub 1)}{
          \fn{Subderiv}((s)_{i+1}, (s)_i) \land
          (n=0 \lor \fn{DischargeLabel}((s)_i) \neq n)}}}.
  \end{multline*}
Cathetan koreksi sumber OLPL-177: anak langsung ing tupla kode derivasi diwiwiti saka indeks siji, amarga indeks nol nyimpen cacah premis; rumus Jawa wis nganggo indeks j tambah siji.
Cathetan panarikan OLPL-178: pratelan lawas yen argumen TeX kang ditutup nyopot cakupan kuantor ora bener. Makro mung nyithak kuantor lan isine; kurung logis sumber wis nyakup kabeh tes jalur. Tata brace nested kang ora ngowahi makna dijaga.
Cathetan koreksi sumber SOL6-F078: label nol ora mbebasake asumsi. Mula jalur label nol ora kena ditolak nalika ngliwati aturan label nol; rumus kang kacithak saiki nampa kasus iki. SOL6-F079 uga mriksa kode jalur sah kang ora kosong.
Cathetan koreksi sumber OLPL-181: wates kanggo kode jalur s kudu ngidinake padha karo kode SubtreeSeq nalika kabeh dhaptar iku dhewe dadi jalur; rumus Jawa wis nganggo wates kang ngidinake kesamaan.
\end{proof}

\begin{prop}
  \ollabel{prop:prf-prim-rec}
  Upamane $\Gamma$ iku himpunan rekursif primitif saka
  !!{sentence}s. Relasi $\Prf[\Gamma](x, y)$ kang nyatakake
  ``$x$ iku kode saka !!a{derivation}~$\delta$ kanggo $!A$
  adhedhasar asumsi kang !!{undischarged} ing~$\Gamma$ lan
  $y$ iku wilangan G\"odel saka~$!A$'' iku rekursif primitif.
\end{prop}

\begin{proof}
  Upamane ``$y \in \Gamma$'' diwenehake dening predikat
  rekursif primitif~$R_\Gamma(y)$. Kita kudu nuduhake yen
  $\Prf[\Gamma](x, y)$, kang bener yen lan mung yen $y$ iku
  wilangan G\"odel saka ukara~$!A$ lan $x$~iku kode saka
  !!{derivation} deduksi natural kanthi !!{formula}
  pungkasan~$!A$ lan kabeh asumsi kang !!{undischarged}
  ana ing~$\Gamma$, iku rekursif primitif.

  Miturut \olref{prop:deriv}, sifat $\fn{Deriv}(x)$, kang bener
  yen lan mung yen $x$ iku wilangan G\"odel saka
  !!{derivation}~$\delta$ kang bener ing deduksi natural,
  iku rekursif primitif. Mula $\Prf[\Gamma](x, y)$ bisa
  ditetepake dening
  \begin{align*}
    \Prf[\Gamma](x, y) \defiff {}
    & \fn{Deriv}(x) \land \fn{EndFmla}(x) = y \land {} \\
    & \bforall{z < x}{(\fn{OpenAssum}(z, x) \lif R_\Gamma(z))}.
  \end{align*}
\end{proof}

\end{document}
